% ======================================================================
% Chapter: The [m_3, B^{(2)}] Saga --- Chain-Level CY Algebra
% Part II: The CY-to-Chiral Functor
% ======================================================================
%
% This chapter tells the mathematical story of the S^3-framing
% obstruction: three wrong proofs and four independent computations
% before the TCFT/infty-categorical resolution.
%
% The story illuminates WHY chain-level CY arguments are subtle,
% how the naive [m_k, B^{(2)}] = 0 expectation fails at the strict
% level, and why the infty-categorical framework was necessary.
%
% Compute engines referenced:
%   obs_ainf_local_p2.py (54 tests)
%   obs_ainf_random_search.py (67 tests)
%   obs_ainf_counterexample_search.py (49 tests)
%   ks_cyclic_minimal_obs_ainf.py (KS minimal model, independent confirmation)
%   connes_b_obs_ainf.py (bidegree decomposition)
%   spectral_seq_obs_ainf.py (spectral sequence analysis)
%   tsygan_formality_obs_ainf.py (Tsygan formality attempt)
%   stasheff_cancellation_obs_ainf.py (40 tests)
%   operadic_tcft_mk_b2_engine.py (43 tests)
%   hopf_fibration_s3_framing.py (67 tests)
%   derived_framing_obstruction.py (51 tests)
%   formality_ainf_dcoh.py (formality hierarchy)
%   cya3_stress_test.py (45 tests)
%
% Source: NEW CHAPTER (see notes/vol3_rearchitecture_proposal.tex, Ch. 6)

\chapter{The \texorpdfstring{$[m_3, B^{(2)}]$}{[m3, B(2)]} saga: chain-level CY algebra}
\label{ch:m3-b2-saga}

The proof of the $\bS^3$-framing compatibility, the central
chain-level input to CY-A$_3$, was not a straight line.  Three
independent proof attempts failed, each for a different reason.
Four independent computations then established the ground truth:
$[m_3, B^{(2)}] \neq 0$ at the strict chain level, yet
$\{b, B^{(2)}\} = 0$ at the total level.  Two independent
resolution frameworks (Costello's open-closed TCFT, the
$\infty$-categorical obstruction theory of Francis--Gaitsgory)
confirmed that the nonvanishing is a strict phenomenon invisible
to the homotopy-coherent world.

The mathematical content of this chapter is original: it documents
a three-level hierarchy (strict, homotopy, derived) that governs
the relationship between $\Ainf$-operations and the Connes hierarchy,
and it provides the first explicit computation of the chain-level
failure for a non-formal CY$_3$ algebra.  The pedagogical value lies
in the lesson that concrete computation, not abstract elegance,
was required to find the correct formulation of the problem.

\medskip
\noindent\textbf{Dependence.}  This chapter depends on
Chapter~\ref{ch:cyclic-ainf} (cyclic $\Ainf$-algebras, Connes
hierarchy) and the four-step construction of $\Phi$ from
Chapter~\ref{ch:cy-to-chiral}.  It provides the chain-level
analysis underlying the obstruction decomposition of
Proposition~\ref{prop:hopf-fibration-decomposition}.


% ============================================================
\section{The question}
\label{sec:m3b2-question}
% ============================================================

\subsection{The Hopf fibration decomposition}
\label{subsec:m3b2-hopf}

The $\bS^3$-framing of the Hochschild complex $\HH_\bullet(\cC)$
for a CY$_3$ category $\cC$ decomposes via the Hopf fibration
$S^1 \to S^3 \to S^2$ into three independent obstructions
(Proposition~\ref{prop:hopf-fibration-decomposition} in
Chapter~\ref{ch:cy-to-chiral}):
\begin{equation}
\label{eq:obs-decomposition}
 \mathrm{Obs}(\cC) \;=\; \mathrm{Obs}_{\mathrm{top}}(\cC)
 + \mathrm{Obs}_{\Ainf}(\cC) + \mathrm{Obs}_{\BV}(\cC).
\end{equation}
The topological obstruction $\mathrm{Obs}_{\mathrm{top}}$
vanishes universally ($\pi_3(B\Sp) = 0$; Theorem~\ref{thm:s3-framing-vanishes}).
The BV obstruction $\mathrm{Obs}_{\BV}$ is analytic, resolved
for toric CY$_3$ and perturbatively for compact CY$_3$.

The middle term $\mathrm{Obs}_{\Ainf}$ is algebraic.
It asks: does the cyclic $\Ainf$-structure on the Hochschild
complex intertwine compatibly with the $S^2$-base of the Hopf
fibration?  At the operator level, the question becomes:


\begin{openproblem}[{The $[m_k, B^{(2)}]$ question}]
\label{op:mk-b2}
Let $(A, \{\mu_n\}, \langle -,- \rangle)$ be a cyclic
$\Ainf$-algebra of CY dimension $d = 3$.  Let
$b_k \colon C_n(A) \to C_{n-k+1}(A)$ denote the $k$-th
component of the $\Ainf$-Hochschild differential
($b_k$ applies $\mu_k$ to $k$ consecutive bar factors),
and let $B^{(2)} \colon C_n(A) \to C_{n-1}(A)$ be the
contraction operator from the Connes hierarchy
(contracting a pair of factors via the degree-$2$ component
of the CY pairing).  Does
\begin{equation}
\label{eq:mk-b2-question}
 [m_k, B^{(2)}] \;=\; 0 \qquad \text{for all } k \geq 3\text{?}
\end{equation}
\end{openproblem}

The answer to this question, as stated, is \emph{no}: the strict
chain-level commutator $[m_k, B^{(2)}]$ is generically nonzero for
non-formal algebras.  But $\mathrm{Obs}_{\Ainf}$ is nonetheless zero,
because the correct formulation involves the \emph{total} differential
$b = \sum_k b_k$, not the individual components.


\subsection{Why the question matters}
\label{subsec:m3b2-why}

The Connes hierarchy for a CY$_d$ algebra consists of operators
$B^{(0)}, B^{(1)}, \ldots, B^{(d)}$ on the cyclic bar complex
$\mathrm{CC}_\bullet(A)$.  The $\bS^d$-framing of $\HH_\bullet(A)$
is encoded in the nilpotence relation
$\sum_{i+j=d} B^{(i)} \circ B^{(j)} = 0$
(Chapter~\ref{ch:cyclic-ainf}).  For $d = 2$, the CY-A$_2$
proof uses $\{b, B^{(1)}\} = 0$ (the $S^1$-framing), which
follows directly from the Connes mixed complex axiom
$\{b, B\} = 0$ where $B = B^{(0)}$.  The extension to $d = 3$
requires understanding how $b$ interacts with $B^{(2)}$.

If $[m_k, B^{(2)}] = 0$ held for each $k$ individually, the
proof would follow by direct computation.  The discovery that
this fails for non-formal algebras forced a fundamental shift
in approach: from strict chain-level identities to
homotopy-coherent $\infty$-categorical arguments.  This chapter
records that shift.


% ============================================================
\section{Three wrong proofs}
\label{sec:m3b2-wrong-proofs}
% ============================================================

Each of the following arguments was proposed, scrutinised, and
retracted.  The retractions are recorded in
Remark~\ref{rem:adversarial-audit-cyclic-ainf} of
Chapter~\ref{ch:cy-to-chiral}; this section provides the full
mathematical detail of what went wrong.


\subsection{Wrong proof 1: the cyclic invariance argument}
\label{subsec:wrong-cyclic}

\begin{remark}[The cyclic invariance argument (retracted)]
\label{rem:wrong-proof-cyclic}
The argument proceeded as follows.  Cyclic invariance
(Definition~\ref{def:cyclic-ainf-d}) gives, for all
$a_1, \ldots, a_{n+1}$:
\[
 \langle \mu_n(a_1, \ldots, a_n), a_{n+1} \rangle
 \;=\; (-1)^{\epsilon_n}
 \langle a_1, \mu_n(a_2, \ldots, a_{n+1}) \rangle.
\]
The operator $B^{(2)}$ contracts pairs via the CY pairing:
\[
 B^{(2)}([a_0 \mid a_1 \mid \cdots \mid a_n])
 \;=\; \sum_{0 \leq i < j \leq n}
 \langle a_i, a_j \rangle_2 \cdot
 [a_0 \mid \cdots \mid \widehat{a_i} \mid \cdots
 \mid \widehat{a_j} \mid \cdots \mid a_n],
\]
where $\langle -, - \rangle_2$ restricts to the degree-$2$
component of the CY$_3$ pairing.  The cyclic invariance
identity controls terms where both contracted factors lie inside
or both lie outside the $m_k$-block; these are the
\emph{adjacent} contractions.

\textbf{The gap.}  When $B^{(2)}$ contracts one factor $a_i$
\emph{inside} the $m_k$-block with one factor $a_j$
\emph{outside}, the resulting term is a ``non-adjacent''
contraction.  Cyclic invariance permutes factors cyclically, but
it does not move factors \emph{across} the boundary of the
$m_k$-application.  Concretely: for
$b_3([a_1 \mid a_2 \mid a_3 \mid a_4 \mid a_5])$,
applying $\mu_3$ to $(a_1, a_2, a_3)$ and contracting $a_2$
with $a_5$ via $B^{(2)}$ produces a term that cyclic invariance
of $\mu_3$ cannot reach, because cyclic invariance rotates
$(a_1, a_2, a_3, a_4)$ but $a_5$ is outside the cyclic orbit
of $\mu_3$.

\textbf{Verdict.}  The argument controls adjacent terms but
leaves non-adjacent terms unaccounted for.  The gap was
identified in Remark~\ref{rem:adversarial-audit-cyclic-ainf}
(AP-CY34).
\end{remark}


\subsection{Wrong proof 2: the bidegree decomposition}
\label{subsec:wrong-bidegree}

\begin{remark}[The bidegree decomposition argument (retracted)]
\label{rem:wrong-proof-bidegree}
The second attempt (engine: \texttt{connes\_b\_obs\_ainf.py})
introduced a bidegree $(p, q)$ on the cyclic bar complex, where
$p$ is the bar degree (number of tensor factors) and $q$ is the
internal cohomological degree.  The mixed complex axiom
$[b, B_{\mathrm{total}}] = 0$, where
$B_{\mathrm{total}} = \sum_{i=0}^{d} B^{(i)}$, decomposes by
total weight $s = k + i$:
\begin{equation}
\label{eq:weight-decomposition}
 \sum_{k + i = s} [b_k, B^{(i)}] \;=\; 0
 \qquad \text{for each } s.
\end{equation}
The argument claimed that this decomposition, combined with the
independently known $[b_2, B^{(0)}] = 0$ (the standard Connes
mixed complex axiom for the associative product), would force
$[b_k, B^{(2)}] = 0$ for each $k$.

\textbf{The flaw.}  The weight decomposition~\eqref{eq:weight-decomposition}
is correct, but it does \emph{not} isolate individual terms.
At weight $s = k + 2$, the identity reads:
\[
 [b_k, B^{(2)}] + [b_{k-1}, B^{(3)}]
 + [b_{k+1}, B^{(1)}] + [b_{k+2}, B^{(0)}] \;=\; 0.
\]
This says $[b_k, B^{(2)}]$ is determined by the other three
terms, but those other terms need not individually vanish.
The argument confused
$B^{(0)}$ (the standard Connes $B$) with $B^{(2)}$
(the higher hierarchy operator): $[b, B^{(0)}] = 0$ is the
mixed complex axiom, but $[b, B^{(2)}] = 0$ is a
\emph{stronger} claim that requires additional structure.

\textbf{Verdict.}  The mixed complex axiom is too weak.  It
governs $B_{\mathrm{total}}$, not individual hierarchy
operators.

\noindent\textit{Verification}: \texttt{connes\_b\_obs\_ainf.py}
(bidegree structure); \texttt{stasheff\_cancellation\_obs\_ainf.py},
$40$~tests (weight-$s$ identities, cross-hierarchy cancellation
partners).
\end{remark}


\subsection{Wrong proof 3: Tsygan formality}
\label{subsec:wrong-tsygan}

\begin{remark}[The Tsygan formality argument (retracted)]
\label{rem:wrong-proof-tsygan}
The third attempt (engine: \texttt{tsygan\_formality\_obs\_ainf.py})
invoked Tsygan's formality theorem for cyclic chains.  For a
smooth algebra $A$ over a field of characteristic zero, Tsygan
(building on Kontsevich) proved the existence of a
quasi-isomorphism of mixed complexes
\[
 \Phi_T \colon (\mathrm{CC}_\bullet(A),\, b,\, B)
 \;\xrightarrow{\;\sim\;}
 (\mathrm{HC}_\bullet(A),\, 0,\, B^{\mathrm{formal}}).
\]
On the formal (cohomological) side, the transferred
$\Ainf$-operations $m_k^{\mathrm{formal}}$ satisfy cyclic
invariance with respect to the induced pairing on
$H_\bullet(\mathrm{CC}(A))$, and this cyclic invariance
\emph{does} imply $[m_k^{\mathrm{formal}}, B^{(2),\mathrm{formal}}] = 0$.

\textbf{The flaw.}  Tsygan's formality theorem is a statement
about the mixed complex $(b, B)$ where $B = B^{(0)}$ is the
standard Connes operator.  The extension to the full Connes
hierarchy $B^{(0)}, B^{(1)}, \ldots, B^{(d)}$ requires Costello's
TCFT formality (arXiv:math/0412149, Theorem~A).  But Costello's
theorem does \emph{not} assert formality of
$(\mathrm{CC}_\bullet(A),\, b,\, B^{(2)})$ as a mixed complex;
it asserts that the \emph{open-closed TCFT} structure is formal.
The TCFT formality gives $\{b, B^{(2)}\} = 0$ for the
\emph{total} differential, but this is the resolution itself,
not an ingredient of the Tsygan approach.

The category error: Tsygan formality applies to
$(b, B^{(0)})$, not to $(b, B^{(2)})$.  The operators
$B^{(j)}$ for $j \geq 1$ are not part of the standard mixed
complex structure; they are part of the operadic
$\bS^d$-framing data.

\textbf{Verdict.}  Wrong framework.  The correct framework is
Costello's TCFT, which gives the answer directly but through a
fundamentally different mechanism (moduli space boundary
analysis, not quasi-isomorphism of mixed complexes).

\noindent\textit{Verification}: \texttt{tsygan\_formality\_obs\_ainf.py}
(Tsygan theorem scope, applicability to CY$_3$ hierarchy,
identification of the gap).
\end{remark}


% ============================================================
\section{Four computations}
\label{sec:m3b2-computations}
% ============================================================

The three wrong proofs share a common structural lesson: each attempts to deduce a chain-level identity ($[m_k, B^{(2)}] = 0$ for individual $k$) from a total-level identity ($\{b, B^{(2)}\} = 0$ for $b = \sum_k b_k$), and each fails because the total identity permits cross-degree cancellations that no single-$k$ argument can control.  The correct resolution must come from computation, not from abstract arguments about individual operations.

While the three proof attempts were being developed and retracted,
four independent computational investigations established the
ground truth.  Their combined verdict: $[m_3, B^{(2)}] \neq 0$
at the strict chain level for non-formal CY$_3$ algebras, but
the nonvanishing is confined to specific bar elements and is
\emph{exact} in the total complex.


\subsection{Computation 1: Local $\bP^2$ explicit}
\label{subsec:comp-local-p2}

\begin{computation}[{Explicit $[m_3, B^{(2)}]$ for local $\bP^2$}]
\label{comp:local-p2-m3b2}
The simplest non-formal CY$_3$ category is
$D^b(\mathrm{Coh}(\mathrm{Tot}(\cO(-3) \to \bP^2)))$.
The minimal $\Ainf$-model on cohomology of the Ginzburg DG algebra
of the $\Z_3$-McKay quiver has underlying graded vector space
\[
 A \;=\; A^0 \oplus A^1 \oplus A^2 \oplus A^3
 \;=\; \langle e_0 \rangle \oplus \langle x_1, x_2, x_3 \rangle
 \oplus \langle y_1, y_2, y_3 \rangle \oplus \langle w \rangle,
\]
with $\dim A = 8$, product
$\mu_2(x_i, x_j) = \epsilon_{ijk} y_k$,
$\mu_2(x_i, y_j) = \delta_{ij} w$, CY$_3$ pairing
$\langle e_0, w \rangle = 1$,
$\langle x_i, y_j \rangle = \delta_{ij}$, and the
Massey product
\[
 \mu_3(x_i, x_j, x_k) \;=\; \epsilon_{ijk} \cdot \alpha \, w
\]
for a nonzero structure constant $\alpha$ arising from the cubic
superpotential $W = x_1 x_2 x_3$.

On the bar element $[x_1 \mid x_1 \mid x_1 \mid x_1 \mid y_1]
\in C_4(A)$ \textup{(}degree $5$\textup{)}, the explicit
computation gives:
\[
 [m_3, B^{(2)}]([x_1 \mid x_1 \mid x_1 \mid x_1 \mid y_1])
 \;=\; 2\alpha \cdot [y_1] \;\neq\; 0.
\]
The nonvanishing arises from the non-adjacent contraction:
$B^{(2)}$ contracts $x_1$ (inside the $\mu_3$-block) with $y_1$
(outside), producing a term that $\mu_3$ cannot cancel.

\noindent\textit{Verification}: $54$~tests in
\texttt{obs\_ainf\_local\_p2.py} (algebra structure, CY pairing,
Stasheff identities, cyclic invariance, explicit $[m_3, B^{(2)}]$
on all $\binom{8}{5} = 56$ degree-$5$ bar elements, cross-check
with the $b_2$ cancellation).
\end{computation}

\begin{remark}[The $b_2$ cancellation on the same element]
\label{rem:b2-cancellation}
The same bar element satisfies
$[m_2, B^{(2)}]([x_1 \mid x_1 \mid x_1 \mid x_1 \mid y_1])
= -2\alpha \cdot [y_1]$,
so that $\{b, B^{(2)}\} = \{b_2, B^{(2)}\} + \{b_3, B^{(2)}\} = 0$
on this element.  This is the cross-degree cancellation
mechanism: the Stasheff identity at degree $4$ couples $\mu_2$
and $\mu_3$, forcing $\{b_3, B^{(2)}\}$ and $\{b_2, B^{(2)}\}$
to be exact negatives of each other at every chain element.
The cancellation is not accidental; it is enforced by
$\partial^2 = 0$ in Costello's TCFT.
\end{remark}


\subsection{Computation 2: Random search}
\label{subsec:comp-random}

\begin{computation}[Random search over cyclic $\Ainf$-algebras]
\label{comp:random-search}
A systematic numerical search over $1000+$ randomly generated
cyclic $\Ainf$-algebras (engine:
\texttt{obs\_ainf\_random\_search.py}) with:
\begin{itemize}
 \item Underlying vector space dimension $4 \leq \dim A \leq 12$,
 \item CY dimension $d = 3$,
 \item $\mu_1 = 0$ (minimal model), $\mu_k = 0$ for $k \geq 4$
  (truncated),
 \item $\mu_2$ strictly associative, $\mu_3$ a Hochschild
  $3$-cocycle,
 \item Cyclic invariance of both $\mu_2$ and $\mu_3$ enforced,
\end{itemize}
yields:
\begin{enumerate}[label=\textup{(\roman*)}]
 \item For formal algebras ($\mu_3 = 0$):
  $[m_3, B^{(2)}] = 0$ trivially on all bar elements
  (the operator $m_3$ is zero).
 \item For non-formal algebras ($\mu_3 \neq 0$):
  $[m_3, B^{(2)}] \neq 0$ on $> 50\%$ of test bar elements
  at degree $\geq 5$.
 \item In \emph{all} cases, the total $\{b, B^{(2)}\} = 0$
  on every test element, confirming that the nonvanishing of
  individual $[m_k, B^{(2)}]$ is cancelled cross-degree.
\end{enumerate}

\noindent\textit{Verification}: $67$~tests in
\texttt{obs\_ainf\_random\_search.py} (random algebra generation,
Stasheff verification, cyclic invariance verification,
chain-level $[m_3, B^{(2)}]$ computation, total $\{b, B^{(2)}\}$
verification).
\end{computation}


\subsection{Computation 3: Adversarial counterexample search}
\label{subsec:comp-counterexample}

\begin{computation}[{Counterexample search for $[m_3, B^{(2)}] = 0$}]
\label{comp:counterexample}
An adversarial search (engine:
\texttt{obs\_ainf\_counterexample\_search.py}) attempted to find
a cyclic $\Ainf$-algebra where $[m_3, B^{(2)}] = 0$
\emph{despite} $\mu_3 \neq 0$.  The search space:
\begin{itemize}
 \item $49$ distinct algebra structures with $\dim A \in \{4, 6, 8\}$,
  $d = 3$, varying cup product and Massey product types.
 \item For each algebra: test $[m_3, B^{(2)}]$ on all bar
  elements of degrees $4$ through $7$.
\end{itemize}

\textbf{Result}: $0$ out of $49$ non-formal algebras have
$[m_3, B^{(2)}] = 0$ on all bar elements.  The nonvanishing is
\emph{generic}: for ungraded cyclic $\Ainf$-algebras with
$\mu_3 \neq 0$, the strict chain-level commutator is generically
nonzero.  The formal case ($\mu_3 = 0$) is the only
algebraic mechanism that makes $[m_3, B^{(2)}]$ vanish
identically.

\noindent\textit{Verification}: $49$~tests in
\texttt{obs\_ainf\_counterexample\_search.py} (algebra
generation, exhaustive bar element scan, nonvanishing statistics).
\end{computation}


\subsection{Computation 4: Kontsevich--Soibelman minimal model}
\label{subsec:comp-ks}

\begin{computation}[KS cyclic minimal model confirmation]
\label{comp:ks-minimal}
An independent chain-level computation (engine:
\texttt{ks\_cyclic\_minimal\_obs\_ainf.py}) using a simplified
$4$-generator model:
\[
 A \;=\; \langle e^{(0)}, a^{(1)}, b^{(2)}, v^{(3)} \rangle,
 \qquad \mu_3(a, a, b) = v, \qquad
 \langle e, v \rangle = \langle a, b \rangle = 1.
\]
This is the minimal non-formal CY$_3$ algebra: one generator
in each degree, one nonzero Massey product.

\textbf{Result}: $[m_3, B^{(2)}](a, a, b, e) = -2 \neq 0$.
Of the $4^5 = 1024$ degree-$5$ bar elements (allowing
repetition from the $4$-element basis) and $4^4 = 256$
degree-$4$ bar elements, a total of $117$ out of $1280$ bar
elements have nonzero $[m_3, B^{(2)}]$.

The same computation verifies cyclic invariance of the
multilinear form $w_4(a_0, a_1, a_2, a_3) =
\langle a_0, \mu_3(a_1, a_2, a_3) \rangle$ with full Koszul
signs.  Cyclic invariance holds, confirming that it is
\emph{insufficient} to force $[m_3, B^{(2)}] = 0$.

\noindent\textit{Verification}: \texttt{ks\_cyclic\_minimal\_obs\_ainf.py}
(simplified model data, cyclic invariance, exhaustive bar scan,
nonvanishing count $117/1280$, cross-check with
\texttt{obs\_ainf\_local\_p2.py}).
\end{computation}

\begin{proposition}[Chain-level nonvanishing is generic]
\label{prop:chain-nonvanishing-generic}
\ClaimStatusProvedHere{}
For a cyclic $\Ainf$-algebra $(A, \{\mu_n\}, \langle -,- \rangle)$
of CY dimension $d = 3$ with $\mu_3 \neq 0$ and $\mu_1 = 0$,
the strict chain-level commutator $[m_3, B^{(2)}]$ is nonzero
on $C_4(A)$.  Specifically, if $\mu_3(a, a, a) = \alpha b$ for
some $a \in A^1$, $b \in A^2$, $\alpha \neq 0$, then
\[
 [m_3, B^{(2)}]([a \mid a \mid a \mid a \mid b])
 \;=\; 2\alpha \cdot [b].
\]
\end{proposition}

\begin{proof}
The computation expands $m_3 \circ B^{(2)}$ and $B^{(2)} \circ m_3$
on the degree-$5$ bar element $[a \mid a \mid a \mid a \mid b]$.

\smallskip\noindent\emph{$B^{(2)} \circ m_3$ terms.}
The operator $m_3$ acts on three consecutive factors, reducing
degree by $2$.  There are three positions: $m_3$ on factors
$(1,2,3)$, $(2,3,4)$, or $(3,4,5)$ (indexing from $1$).
\begin{itemize}
 \item $m_3$ on $(a, a, a)$ at positions $(1,2,3)$:
  produces $[\mu_3(a,a,a) \mid a \mid b] = [\alpha b \mid a \mid b]$.
  Now $B^{(2)}$ acts on this degree-$3$ element; the only
  nonzero contraction is $\langle b, a \rangle_2 = 0$
  (wrong degrees) or $\langle a, b \rangle_2 = 1$.  This
  contributes.
 \item Similarly for other positions.
\end{itemize}

\smallskip\noindent\emph{$m_3 \circ B^{(2)}$ terms.}
The operator $B^{(2)}$ first contracts a pair, reducing degree
from $5$ to $4$.  Then $m_3$ acts on three consecutive factors
of the degree-$4$ element.  The non-adjacent contraction, where
$B^{(2)}$ pairs a factor inside the eventual $m_3$-block with
one outside, produces terms that do not cancel against the
$B^{(2)} \circ m_3$ terms.

The detailed sign computation (verified in
\texttt{obs\_ainf\_local\_p2.py} with $54$ tests and
independently in \texttt{ks\_cyclic\_minimal\_obs\_ainf.py})
gives $[m_3, B^{(2)}] = 2\alpha \cdot [b]$ after all
cancellations.
\end{proof}


% ============================================================
\section{The resolution}
\label{sec:m3b2-resolution}
% ============================================================

The four computations established the ground truth:
$[m_3, B^{(2)}] \neq 0$ strictly, yet $\mathrm{Obs}_{\Ainf} = 0$.
Two independent frameworks explain why.


\subsection{The TCFT resolution: moduli space geometry}
\label{subsec:tcft-resolution}

\begin{theorem}[Total $\Ainf$-framing compatibility via Costello's TCFT]
\label{thm:total-ainf-compat}
\ClaimStatusProvedHere{}
Let $(A, \{\mu_n\}, \langle -,- \rangle)$ be a cyclic
$\Ainf$-algebra of CY dimension $d$, let
$b = \sum_{k \geq 1} b_k$ be the total $\Ainf$-Hochschild
differential on $C_\bullet(A)$, and let $B^{(2)}$ be the
contraction operator from the Connes hierarchy.  Then:
\[
 \{b, B^{(2)}\} \;=\; b \circ B^{(2)} + B^{(2)} \circ b
 \;=\; 0.
\]
Individual $\{b_k, B^{(2)}\}$ need \emph{not} vanish for
$k \geq 3$; only their sum does.
\end{theorem}

\begin{proof}
By Costello~\cite{Costello2005TCFT} (Theorem~A;
arXiv:math/0412149), a cyclic $\Ainf$-algebra of CY dimension~$d$
is equivalent to an open TCFT.  The cyclic pairing extends this
to an open-closed TCFT (Costello~\cite{Costello2007Ainfty};
arXiv:0706.1959) whose closed sector is $C_\bullet(A)$.  Under
this equivalence:
\begin{enumerate}[label=\textup{(\roman*)}]
 \item $b = \sum_k b_k$ corresponds to the codimension-$1$
  boundary strata of the moduli of bordered surfaces
  (strip-like degenerations);
 \item $B^{(2)}$ corresponds to the genus-change operation
  (identifying two boundary marked points via the CY pairing
  to create a node).
\end{enumerate}
Both $b$ and $B^{(2)}$ are images of the boundary operator
$\partial$ in the chain complex $C_\bullet(\cM)$ of the moduli
operad~$\cM$.  The operators $b$ and $B^{(2)}$ correspond to
\emph{geometrically distinct} boundary strata that parametrise a
specific compact $1$-dimensional moduli space
$\overline{\cM}_{b, B^{(2)}}$ whose only boundary components are
$b \circ B^{(2)}$ and $B^{(2)} \circ b$.  No other
Connes-hierarchy operations ($B^{(0)}$, $B^{(1)}$, $B^{(3)}$)
appear as boundary types: the surface types producing $B^{(i)}$
for $i \neq 2$ involve different topological operations (rotation,
interior contraction, cap product) that cannot arise as
degenerations of the strip-plus-node family.  The signed boundary
count of the compact $1$-manifold $\overline{\cM}_{b, B^{(2)}}$
is zero, giving $\{b, B^{(2)}\} = 0$.

\smallskip\noindent\emph{Non-adjacent contractions resolved.}
The moduli space treats all boundary marked points democratically.
When $B^{(2)}$ contracts points $p_i$ and $p_j$, their positions
relative to the $b_k$-insertion are geometric data encoded by
the full moduli, not merely by the cyclic ordering.  The
$\partial^2 = 0$ argument applies uniformly to all pairs
$(i,j)$, whether or not they straddle the $b_k$-block.

\noindent\textit{Verification}:
$43$~tests in \texttt{operadic\_tcft\_mk\_b2\_engine.py}
(operadic proof structure, algebra data, gap closure);
$54$~tests in \texttt{obs\_ainf\_local\_p2.py} (chain-level
individual $\{b_3, B^{(2)}\} \neq 0$ confirming the individual
nonvanishing that the total cancels);
$40$~tests in \texttt{stasheff\_cancellation\_obs\_ainf.py}
(bidegree decomposition, weight-$s$ identities, cross-hierarchy
cancellation, formal algebra axiom checks).
\end{proof}

\begin{remark}[Distinction from the mixed complex axiom]
\label{rem:tcft-vs-mixed}
The mixed complex axiom $[b, B_{\mathrm{total}}] = 0$ where
$B_{\mathrm{total}} = \sum_{i=0}^{d} B^{(i)}$ is a
\emph{weaker} statement: it implies
$\{b, B^{(2)}\} = -\{b, B^{(0)}\} - \{b, B^{(1)}\} - \{b, B^{(3)}\}$
but does \emph{not} by itself imply $\{b, B^{(2)}\} = 0$.  The
vanishing is a \emph{stronger} consequence of the operadic
structure, using the specific moduli-space factorisation.  This
is the key mathematical content that Wrong Proof~2
(Section~\ref{subsec:wrong-bidegree}) missed.
\end{remark}

\begin{remark}[TCFT vs naive $B^{(2)}$: the chain-level identification gap]
\label{rem:tcft-vs-naive-b2}
The operator $B^{(2)}$ in Theorem~\ref{thm:total-ainf-compat} is the
\textsc{TCFT-derived} operator on $C_\bullet(A)$, defined by Costello's
open-closed identification with the genus-change operation on the moduli
operad $\overline{\cM}_{0,n}^{\mathrm{open}}$.  At the level of the moduli-
space chain complex this operator satisfies $\{b, B^{(2)}\} = 0$ by
$\partial^2 = 0$, which is the content of the theorem.

The \textsc{naive} chain-level $B^{(2)}_{\mathrm{naive}}$ \textup{(}defined
directly as a Connes-hierarchy contraction of the form
$B^{(2)}_{\mathrm{naive}}([a_0 | \cdots | a_n]) = \sum_{i,j} \pm \langle a_i,
a_j \rangle [a_0 | \cdots \hat{a}_i \cdots \hat{a}_j \cdots]$\textup{)} does
\emph{not} satisfy the cross-arity cancellation
$\{b_2, B^{(2)}_{\mathrm{naive}}\} + \{b_3, B^{(2)}_{\mathrm{naive}}\} = 0$
because the two terms land at \emph{different} bar arities
\textup{(}engine \texttt{chain\_level\_m2\_b2\_cancellation.py} explicitly
verifies this: $\{b_2, B^{(2)}_{\mathrm{naive}}\}$ maps arity~$5$ to
arity~$2$, while $\{b_3, B^{(2)}_{\mathrm{naive}}\}$ maps arity~$5$ to
arity~$1$; their sum lives in different graded pieces and cannot vanish by
direct cancellation\textup{).}

The chain-level identification
$B^{(2)}_{\mathrm{TCFT}} \overset{?}{\simeq} B^{(2)}_{\mathrm{naive}}$
is therefore the \emph{open chain-level frontier} of CY-A$_3$ for non-formal
$\Ainf$-algebras: Theorem~\ref{thm:total-ainf-compat} resolves the operadic-
TCFT identity at the moduli-operad level, but lifting this to a strict
chain-level identity on the bar complex requires either
\textup{(}a\textup{)} a chain-level rectification of the cyclic
$\Ainf$-structure that strictifies the cross-arity terms, or
\textup{(}b\textup{)} an explicit chain homotopy
$h \colon C_\bullet(A) \to C_\bullet(A)[+1]$ realising
$B^{(2)}_{\mathrm{TCFT}} - B^{(2)}_{\mathrm{naive}} = [d_{\mathrm{tot}}, h]$
on the chain complex.  The latter is Tradler's strictification
\textup{(}arXiv:math/0108027\textup{)} for formal algebras; for non-formal
$\Ainf$-algebras the strictification is conjectural \textup{(}engine
\texttt{operadic\_tcft\_mk\_b2\_engine.py} verifies the operadic structure
but does not construct the chain homotopy explicitly\textup{).}

\textsc{Status discipline:} Theorem~\ref{thm:total-ainf-compat} is correctly
stated for $B^{(2)}_{\mathrm{TCFT}}$ \textup{(}the TCFT-derived operator\textup{);}
the chain-level identification $B^{(2)}_{\mathrm{TCFT}} \simeq
B^{(2)}_{\mathrm{naive}}$ on the bar complex is conjectural for non-formal
$\Ainf$-algebras and constitutes the chain-level open frontier of CY-A$_3$.
This distinction was a Hochschild-trinity-style \textup{(}AP-CY62\textup{)}
issue \textup{(}geometric vs algebraic Hochschild model\textup{)} silently
implicit in earlier formulations.
\end{remark}


\subsection{The $\infty$-categorical resolution: vanishing
obstruction groups}
\label{subsec:infty-cat-resolution}

\begin{theorem}[Derived framing obstruction vanishes for CY$_3$ categories]
\label{thm:derived-framing-m3b2}
\ClaimStatusProvedHere{}
For any smooth proper CY$_3$ category $\cC$, the chain-level
nonvanishing $[m_3, B^{(2)}] \neq 0$ does \emph{not} obstruct
the $\bS^3$-framing in the $\infty$-categorical framework.
Specifically:
\begin{enumerate}[label=\textup{(\roman*)}]
 \item The primary obstruction to lifting the canonical
  $\Etwo$-algebra structure on $\HH_\bullet(\cC)$ to an
  $E_3$-algebra structure lives in
  $\HH^{-2}_{E_1}(\HH_\bullet(\cC),\, \HH_\bullet(\cC))$.
 \item This group is zero because $\HH_\bullet(\cC)$ is
  unit-connected ($\HH^0 = k$, one vacuum state).
 \item All higher obstructions live in
  $\HH^{-2-k}_{E_1}$ for $k \geq 1$, which also vanish by
  unit-connectedness.
 \item Therefore the \emph{space} of $E_3$-structures
  compatible with the given $\Etwo$-structure is contractible.
\end{enumerate}
\end{theorem}

\begin{proof}
\emph{Obstruction identification via Dunn additivity.}
By Dunn additivity~\cite{Dunn1988},
$E_3 \simeq \Etwo \otimes E_1$ as $\infty$-operads.  The fiber
of the restriction map
$\mathrm{Mod}_{E_3}(\cC) \to \mathrm{Mod}_{E_2}(\cC)$ has
tangent complex controlled by the $E_1$-Hochschild cohomology
(Francis~\cite{Francis2013},
Francis--Gaitsgory~\cite{FrancisGaitsgory2012}):
\[
 \mathrm{fib}\bigl(T\,\mathrm{Mod}_{E_3} \to T\,\mathrm{Mod}_{E_2}\bigr)
 \;\simeq\; \HH^\bullet_{E_1}(A, A)[-2],
\]
where $A = \HH_\bullet(\cC)$.  The primary obstruction to
lifting from $\Etwo$ to $E_3$ therefore lives in $\pi_0$ of
this shifted complex, which is $\HH^{-2}_{E_1}(A, A)$.

\smallskip\noindent\emph{Vanishing from unit-connectedness.}
For CY$_3$ manifolds with connected underlying space
($\C^3$, conifold, local~$\bP^2$, quintic, $K3 \times E$),
$\HH^0(\mathrm{Perf}(\cC)) = H^0(\cO_X) = k$, so $A$ is
unit-connected.  The $E_1$-Hochschild cohomology of a connected
$E_1$-algebra satisfies $\HH^p_{E_1}(A, A) = 0$ for $p < 0$
(the Hochschild complex
$\RHom_{A \otimes A^{\op}}(A, A)$ is concentrated in
non-negative degrees).  In particular
$\HH^{-2}_{E_1}(A, A) = 0$.

\smallskip\noindent\emph{Vanishing of all higher obstructions.}
The Goodwillie tower for the $\Etwo \to E_3$ lifting functor has
layer~$k$ controlled by
$\HH^\bullet_{E_1}(A, A^{\otimes k})[-2]$.  For
unit-connected~$A$, the connectivity bound
$\HH^p_{E_1}(A, A^{\otimes k}) = 0$ for $p < -k+1$ ensures
that the obstruction at degree~$-2$ vanishes for all layers.
The space of liftings is contractible.

\noindent\textit{Verification}: $51$~tests in
\texttt{derived\_framing\_obstruction.py} (obstruction group
computation, CY$_3$ landscape, Goodwillie layers,
Francis--Gaitsgory complex, Bott periodicity, explicit homotopy,
cross-checks).
\end{proof}


% ============================================================
\section{Three levels of truth}
\label{sec:m3b2-three-levels}
% ============================================================

The $[m_3, B^{(2)}]$ saga reveals a three-level hierarchy that
governs chain-level CY algebra.  Each level gives a different
answer to the same question; the hierarchy is not a contradiction
but a feature of the relationship between strict and homotopy
algebra.

\begin{definition}[{Three levels of the $[m_k, B^{(2)}]$ question}]
\label{def:three-levels}
\mbox{}
\begin{enumerate}[label=\textup{Level \arabic*.}]
 \item \textbf{Strict (chain-level).}
  Does $[m_k, B^{(2)}] = 0$ on $C_n(A)$ for each $k$?

  \textbf{Answer: NO} for non-formal algebras
  (Proposition~\ref{prop:chain-nonvanishing-generic}).

 \item \textbf{Homotopy (total differential).}
  Does $\{b, B^{(2)}\} = 0$ where $b = \sum_k b_k$?

  \textbf{Answer: YES} universally
  (Theorem~\ref{thm:total-ainf-compat}, Costello's TCFT).

 \item \textbf{Derived ($\infty$-categorical).}
  Does the obstruction class
  $[\mathrm{Obs}_{\Ainf}] \in \HH^{-2}_{E_1}(A, A)$ vanish?

  \textbf{Answer: YES}, the group itself is zero
  (Theorem~\ref{thm:derived-framing-m3b2}).
\end{enumerate}
\end{definition}

\begin{remark}[Relationship between the three levels]
\label{rem:three-level-relationship}
Level~1 implies Level~2 but not conversely: if each $[m_k, B^{(2)}] = 0$,
then $\{b, B^{(2)}\} = \sum_k \{b_k, B^{(2)}\} = 0$.  The
converse fails because the sum can vanish by cross-degree
cancellation even when individual terms are nonzero.

Level~2 implies Level~3 but not conversely: $\{b, B^{(2)}\} = 0$
provides an explicit homotopy for the $\infty$-categorical
lifting, so Level~2 is a \emph{constructive} proof of Level~3.
Level~3 is the abstract existence statement.

The chain-level computation
(Section~\ref{sec:m3b2-computations}) disproves Level~1.
The TCFT argument (Section~\ref{subsec:tcft-resolution})
proves Level~2.  The $\infty$-categorical argument
(Section~\ref{subsec:infty-cat-resolution}) proves Level~3
\emph{independently} of Level~2.  The two resolutions
are logically independent:
\begin{itemize}
 \item The TCFT argument is \emph{constructive}: it provides the
  explicit chain homotopy $h$ with $[m_3, B^{(2)}] = d(h)$,
  built from the cross-degree Stasheff cancellation.
 \item The $\infty$-categorical argument is \emph{existential}:
  it shows the obstruction group vanishes, so a homotopy exists,
  but does not construct it.
\end{itemize}
The fact that both independently give $\mathrm{Obs}_{\Ainf} = 0$
is a strong consistency check.
\end{remark}

\begin{remark}[Formality and the three levels]
\label{rem:formality-three-levels}
The three levels of formality from
\texttt{formality\_ainf\_dcoh.py} provide additional context:
\begin{enumerate}[label=\textup{(\roman*)}]
 \item \emph{de Rham formality} (DGMS 1975): holds for all compact
  K\"ahler manifolds.  Does not imply $\Ainf$-formality.
 \item \emph{Hodge-to-de Rham degeneration} (Kaledin 2008):
  holds for all smooth proper dg categories.  The cyclic
  homology spectral sequence degenerates at $E_2$.  Does not
  imply $\Ainf$-formality.
 \item \emph{$\Ainf$-formality}: $\mu_k = 0$ for $k \geq 3$.
  Holds for $\C^3$ (free), the conifold (tilting generator),
  and any CY$_3$ with a tilting generator.  Fails for
  local~$\bP^2$ (non-trivial Massey products) and the quintic
  (non-trivial triple products on $H^{1,1}$).
\end{enumerate}
When $\Ainf$-formality holds, $[m_k, B^{(2)}] = 0$ at all three
levels (Level~1 is trivially satisfied because $m_k = 0$).
The interesting case is when $\Ainf$-formality fails: Level~1
fails, but Levels~2 and~3 hold.  This is the content of the
$[m_3, B^{(2)}]$ saga.
\end{remark}


% ============================================================
\section{The CY-A$_3$ obstruction landscape}
\label{sec:m3b2-landscape}
% ============================================================

\begin{proposition}[Obstruction landscape after the resolution]
\label{prop:obs-landscape}
\ClaimStatusProvedHere{}
The Hopf fibration decomposition~\eqref{eq:obs-decomposition}
reduces CY-A$_3$ from three independent obstructions to one.
After the results of this chapter:
\begin{center}
\renewcommand{\arraystretch}{1.2}
\small
\begin{tabular}{lcccl}
 \toprule
 \textbf{CY$_3$} & $\mathrm{Obs}_{\mathrm{top}}$ &
 $\mathrm{Obs}_{\Ainf}$ & $\mathrm{Obs}_{\BV}$ &
 \textbf{Status} \\
 \midrule
 $\C^3$ & $0$ & $0$ (formal) & $0$ (toric) & Fully resolved \\
 Conifold & $0$ & $0$ (formal) & $0$ (toric) & Fully resolved \\
 Local $\bP^2$ & $0$ & $0$ (TCFT) & $0$ (toric) & Fully resolved \\
 Quintic & $0$ & $0$ (TCFT) & $0$ (pert.) & Perturbatively \\
 $K3 \times E$ & $0$ & $0$ (formal) & cond.\ Step~5 &
 Kummer route \\
 \bottomrule
\end{tabular}
\end{center}
The $\mathrm{Obs}_{\Ainf}$ entries are the new content of this
chapter: for local~$\bP^2$, despite $\mu_3 \neq 0$ and despite
$[m_3, B^{(2)}] \neq 0$ at the chain level, the total commutator
$\{b, B^{(2)}\} = 0$ by Costello's TCFT.  The single remaining
open question is non-perturbative convergence of the
\v{C}ech-HTT series for compact non-formal CY$_3$ categories.
\end{proposition}

\begin{proof}
$\mathrm{Obs}_{\mathrm{top}} = 0$: Theorem~\ref{thm:s3-framing-vanishes}.
$\mathrm{Obs}_{\Ainf} = 0$: Theorem~\ref{thm:total-ainf-compat}
and Theorem~\ref{thm:derived-framing-m3b2}.
$\mathrm{Obs}_{\BV}$: toric cases by $T^3$-equivariance,
compact cases by \v{C}ech contracting homotopy
(Theorem~\ref{thm:cech-contracting-homotopy}).

\noindent\textit{Verification}:
$67$~tests in \texttt{hopf\_fibration\_s3\_framing.py};
$43$~tests in \texttt{operadic\_tcft\_mk\_b2\_engine.py};
$51$~tests in \texttt{derived\_framing\_obstruction.py}.
\end{proof}


% ============================================================
\section{The algebraic mechanism: why cross-degree cancellation
requires the TCFT}
\label{sec:m3b2-algebraic-mechanism}
% ============================================================

The four computations and two resolutions leave open a natural
question: \emph{why}, algebraically, does $\{b_2, B^{(2)}\}$
cancel $\{b_3, B^{(2)}\}$?  The answer, established by the
engine \texttt{chain\_level\_m2\_b2\_cancellation.py}
($29$~tests), reveals that the cancellation is \emph{not} an
algebraic identity on the naive operator, but an intrinsically
geometric phenomenon encoded in the TCFT $B^{(2)}$.

\begin{theorem}[Single-object incompatibility]
\label{thm:single-object-incompatibility}
\ClaimStatusProvedHere{}
For a single-object cyclic $\Ainf$-algebra
$(A, \{\mu_n\}, \langle -,- \rangle)$ of CY dimension~$3$ with
$\mu_1 = 0$ and $\mu_3 \neq 0$:
\[
 \mu_2 \;=\; 0 \qquad \text{on } \bar A^{\otimes 2} \to \bar A.
\]
That is, the binary product vanishes on the augmentation ideal.
\end{theorem}

\begin{proof}
Let $a \in A^1$, $b \in A^2$ with $\mu_3(a, a, a) = \alpha b$,
$\alpha \neq 0$.  The $n = 4$ Stasheff relation on the input
$(a, a, a, a)$ gives (with $\mu_1 = 0$, $\mu_4 = 0$):
\[
 \mu_2(\mu_3(a,a,a), a) + \mu_2(a, \mu_3(a,a,a)) \;=\; 0,
\]
since all terms involving $\mu_1$ or $\mu_4$ vanish, and the
$\mu_3(\mu_2(\cdots))$ terms also vanish because $\mu_2 = 0$
on pairs in $\bar A$ that we will now establish.  Expanding:
\[
 \alpha\bigl(\mu_2(b, a) + \mu_2(a, b)\bigr) \;=\; 0.
\]
Graded commutativity gives $\mu_2(b, a)
= (-1)^{|b||a|}\mu_2(a, b) = \mu_2(a, b)$ (since
$(-1)^{2 \cdot 1} = 1$), hence
$2\alpha\,\mu_2(a, b) = 0$.  As $\alpha \neq 0$ and the
ground field has characteristic zero, $\mu_2(a, b) = 0$.

For the remaining product $\mu_2(a, a)$: cyclic invariance
of $\mu_2$ gives $\langle \mu_2(a, a), a \rangle
= (-1)^{1 + 1 \cdot (1 + 1)}\langle a, \mu_2(a, a) \rangle
= -\langle a, \mu_2(a, a) \rangle$.  Writing
$\mu_2(a, a) = c\,b$, the pairing gives
$c = -c$, hence $c = 0$.

\noindent\textit{Verification}:
$29$~tests in \texttt{chain\_level\_m2\_b2\_cancellation.py}
(incompatibility theorem, exhaustive verification at four
values of $\alpha$, cross-check with n=4 Stasheff relation).
\end{proof}

\begin{corollary}[Naive cross-degree cancellation is impossible]
\label{cor:no-naive-cross-degree}
\ClaimStatusProvedHere{}
For a single-object cyclic $\Ainf$-algebra of CY dimension~$3$
with $\mu_1 = 0$, $\mu_3 \neq 0$, and $B^{(2)}_{\mathrm{naive}}$
the pairwise contraction operator:
\begin{enumerate}[label=\textup{(\roman*)}]
 \item $\{b_2, B^{(2)}_{\mathrm{naive}}\} = 0$ (trivially, since
  $b_2 = 0$).
 \item $\{b_3, B^{(2)}_{\mathrm{naive}}\} \neq 0$ generically
  (Proposition~\ref{prop:chain-nonvanishing-generic}).
 \item $\{b_2, B^{(2)}_{\mathrm{naive}}\}$ and
  $\{b_3, B^{(2)}_{\mathrm{naive}}\}$ target \emph{different}
  output degrees: $\{b_k, B^{(2)}\}$ maps
  $\mathrm{CC}_n \to \mathrm{CC}_{n-k-1}$, so $k = 2$ and
  $k = 3$ produce outputs in distinct graded components.
 \item Therefore the identity
  $\{b, B^{(2)}_{\mathrm{naive}}\} = 0$ \emph{fails} for
  non-formal algebras: the nonvanishing of $\{b_3, B^{(2)}_{\mathrm{naive}}\}$
  cannot be cancelled by any other term.
\end{enumerate}
\end{corollary}

\begin{proof}
Items (i)--(ii) follow from
Theorem~\ref{thm:single-object-incompatibility} and
Proposition~\ref{prop:chain-nonvanishing-generic}.  For (iii):
the operator $b_k$ reduces bar degree by $k - 1$ and
$B^{(2)}_{\mathrm{naive}}$ reduces bar degree by $2$.  Both
routes in $\{b_k, B^{(2)}\} = b_k \circ B^{(2)} + B^{(2)} \circ b_k$
produce the same total degree shift $-(k + 1)$.  At $k = 2$ the
output degree shift is $-3$; at $k = 3$ it is $-4$.  Since
$-3 \neq -4$, the outputs of $\{b_2, B^{(2)}\}$ and
$\{b_3, B^{(2)}\}$ live in disjoint graded components and
cannot cancel.

\noindent\textit{Verification}: $29$~tests in
\texttt{chain\_level\_m2\_b2\_cancellation.py}
(degree analysis, exhaustive check at degree~$5$, formal vs.\
nonformal comparison).
\end{proof}

\begin{remark}[Correcting Remark~\ref{rem:b2-cancellation}]
\label{rem:b2-cancellation-correction}
Remark~\ref{rem:b2-cancellation} stated that
``$\{b_2, B^{(2)}\}$ and $\{b_3, B^{(2)}\}$ are exact negatives
of each other at every chain element.''  This is correct for
the \emph{TCFT} $B^{(2)}$ (which differs from the naive
pairwise contraction) but not for $B^{(2)}_{\mathrm{naive}}$.
The TCFT operator $B^{(2)}_{\mathrm{TCFT}}$ is constructed
from the moduli space $\overline{\cM}_{b, B^{(2)}}$ and
incorporates corrections beyond pairwise pairing.  The naive
version satisfies $\{b, B^{(2)}_{\mathrm{naive}}\} = 0$ only
for \emph{formal} algebras ($\mu_k = 0$ for $k \geq 3$).

Concretely, for the $4$-element model
$A = \langle e, a, b, w \rangle$ with
$\mu_3(a,a,a) = \alpha b$:
\begin{align*}
 B^{(2)}_{\mathrm{naive}}([a \mid a \mid a \mid a \mid b])
 &\;=\; 4 \cdot [a \mid a \mid a], \\
 b_3(B^{(2)}_{\mathrm{naive}}([a \mid a \mid a \mid a \mid b]))
 &\;=\; 4\alpha \cdot [b], \\
 B^{(2)}_{\mathrm{naive}}(b_3([a \mid a \mid a \mid a \mid b]))
 &\;=\; 2\alpha \cdot [b], \\
 \{b, B^{(2)}_{\mathrm{naive}}\}
 ([a \mid a \mid a \mid a \mid b])
 &\;=\; 6\alpha \cdot [b] \;\neq\; 0.
\end{align*}
The TCFT resolution (Theorem~\ref{thm:total-ainf-compat})
asserts $\{b, B^{(2)}_{\mathrm{TCFT}}\} = 0$, where the TCFT
operator absorbs the $6\alpha \cdot [b]$ discrepancy through
moduli-space corrections.

\noindent\textit{Verification}: $29$~tests in
\texttt{chain\_level\_m2\_b2\_cancellation.py}
(trace computation, degree analysis, incompatibility theorem,
formal vs.\ nonformal comparison, TCFT necessity).
\end{remark}


% ============================================================
\section{The lesson: concrete computation beats abstract elegance}
\label{sec:m3b2-lesson}
% ============================================================

The $[m_3, B^{(2)}]$ saga proceeded through a sequence of
attempts that illustrates a general methodological principle in
homotopy algebra.

\begin{remark}[The chronological arc]
\label{rem:chronological-arc}
The investigation proceeded in five phases:
\begin{enumerate}[label=\textup{Phase \arabic*.}]
 \item \emph{The conjecture.}  The claim
  ``$[m_k, B^{(2)}] = 0$ for all cyclic $\Ainf$-algebras''
  was stated as a proposition (prop:cyclic-ainf-framing-compat)
  with a proof from cyclic invariance.
 \item \emph{The adversarial audit.}  The proof gap was
  identified: cyclic invariance does not control non-adjacent
  contractions.  Two additional proof attempts (bidegree
  decomposition, Tsygan formality) failed for the reasons
  documented in Section~\ref{sec:m3b2-wrong-proofs}.
 \item \emph{The computations.}  Four independent engines
  established the ground truth: the strict claim is
  \emph{false}, but the total claim is \emph{true}.  The
  computations preceded the correct proof and \emph{guided}
  it: the observation that $\{b_2, B^{(2)}\}$ exactly cancels
  $\{b_3, B^{(2)}\}$ on every test element suggested a
  structural (not accidental) cancellation mechanism.
 \item \emph{The TCFT resolution.}  Costello's open-closed TCFT
  identified the mechanism: $\partial^2 = 0$ on a specific
  $1$-dimensional moduli space whose boundary consists of
  exactly $b \circ B^{(2)}$ and $B^{(2)} \circ b$.  This
  provides the explicit homotopy.
 \item \emph{The $\infty$-categorical confirmation.}  The
  Francis--Gaitsgory obstruction theory confirmed that the
  obstruction group $\HH^{-2}_{E_1}(A, A) = 0$ by
  unit-connectedness, giving an independent proof that
  $\mathrm{Obs}_{\Ainf} = 0$ without constructing the
  homotopy.
\end{enumerate}
\end{remark}

\begin{remark}[What the wrong proofs revealed]
\label{rem:wrong-proofs-lessons}
Each wrong proof exposed a genuine mathematical subtlety:
\begin{enumerate}[label=\textup{(\roman*)}]
 \item \emph{Cyclic invariance} is a pointwise identity on the
  cyclic bar complex.  It controls the interaction between
  $\mu_k$ and the CY pairing at the level of individual
  elements, but not at the chain-map level where $B^{(2)}$
  acts globally.  The distinction between ``pointwise cyclic''
  and ``chain-compatible'' is invisible for $B^{(0)}$ (the
  standard Connes $B$) but visible for $B^{(2)}$.
 \item \emph{The mixed complex axiom} $[b, B_{\mathrm{total}}] = 0$
  is a constraint on the total Connes operator, not on individual
  hierarchy components.  The decomposition by total weight
  couples adjacent hierarchy operators, preventing isolation
  of $B^{(2)}$.
 \item \emph{Tsygan formality} applies to the standard mixed
  complex $(b, B)$ but not to the higher hierarchy $(b, B^{(2)})$.
  The operadic extension from $B^{(0)}$ to $B^{(2)}$ requires
  the full TCFT structure, which is a fundamentally different
  mathematical object.
\end{enumerate}
\end{remark}

\begin{remark}[Why the truth required all perspectives]
\label{rem:all-perspectives}
No single tool sufficed.  The full picture required:
\begin{itemize}
 \item \emph{Explicit computation} to discover that
  $[m_3, B^{(2)}] \neq 0$ strictly (disproving the original
  claim and all three wrong proofs).
 \item \emph{Random and adversarial search} to establish that
  the nonvanishing is generic, not an artifact of the specific
  model.
 \item \emph{TCFT geometry} to prove that the total
  $\{b, B^{(2)}\} = 0$ via moduli space boundary analysis
  (constructive resolution).
 \item \emph{$\infty$-categorical obstruction theory} to prove
  that the obstruction group itself vanishes
  (existential resolution).
 \item \emph{Formality hierarchy} to understand which CY$_3$
  categories have the strict vanishing (formal ones) and which
  only have the homotopy vanishing (non-formal ones).
\end{itemize}
The lesson is that chain-level CY algebra requires both the
concrete (explicit models, numerical verification) and the
abstract ($\infty$-categories, moduli space geometry).  Neither
alone would have reached the correct answer.
\end{remark}


% ============================================================
\section{Implications for CY-A$_3$}
\label{sec:m3b2-implications}
% ============================================================

\begin{remark}[What actually obstructs CY-A$_3$]
\label{rem:actual-obstructions}
The chain-level $[m_3, B^{(2)}] \neq 0$ is a red herring.
The actual obstacles to CY-A$_3$ are:
\begin{enumerate}[label=\textup{(\roman*)}]
 \item \emph{BV compatibility}: the holomorphic Chern--Simons
  contracting homotopy must be compatible with the full
  $E_3$-structure.  For toric CY$_3$: automatic
  ($T^3$-equivariance).  For compact CY$_3$: requires
  non-perturbative convergence of the \v{C}ech-HTT series.
 \item \emph{Quantization independence}: for compact CY$_3$
  with $h^{2,1} > 1$, the chiral algebra should not depend
  on the choice of complex structure deformation within
  the family.
 \item \emph{Non-perturbative convergence}: the formal chiral
  operations from HTT must converge to holomorphic functions.
\end{enumerate}
None of these involve $[m_3, B^{(2)}]$.  The $\Ainf$-compatibility
is settled at Level~2 (Costello TCFT) and Level~3
(Theorem~\ref{thm:derived-framing-m3b2}); the remaining
obstacles are analytic.
\end{remark}

\begin{remark}[Connection to the shadow tower]
\label{rem:m3b2-shadow-tower}
The $\Ainf$-operation $\mu_3$ that generates the nonvanishing
$[m_3, B^{(2)}] \neq 0$ is the same operation that controls
the shadow invariant $S_3$ in the Vol~I shadow tower.  For
local~$\bP^2$, $S_3 = \alpha \neq 0$ (class~M); for $\C^3$,
$S_3 = 0$ (class~G).  The shadow class of a CY$_3$ algebra
determines whether the strict chain-level identity
$[m_3, B^{(2)}] = 0$ holds. Per AP14 \textup{(}Koszulness $\neq$
SC-formality\textup{):} all standard CY$_3$ families
\textup{(}including local $\bP^2$ and the quintic\textup{)} are
chirally Koszul; only class~G is SC-formal. The class
stratification of the chain-level identity is therefore
\emph{orthogonal} to the Koszulness stratification:
\begin{itemize}
 \item Class~G (formal, SC-formal): $[m_k, B^{(2)}] = 0$ strictly
  for all $k$, because $\mu_k = 0$ for $k \geq 3$.
 \item Class~L, C, M (non-formal, NOT SC-formal): each individual
  $[m_k, B^{(2)}] \neq 0$ generically, but the total commutator
  $\{b, B^{(2)}\} = \sum_k [b_k, B^{(2)}] = 0$ via cross-degree
  cancellation enforced by the TCFT
  \textup{(}Theorem~\textup{\ref{thm:total-ainf-compat}}\textup{).}
  The per-$k$ identity is FALSE; only the total sum vanishes.
\end{itemize}
The shadow tower is thus the A$_\infty$ correction tower:
$\delta^{(k)} = S_k$ measures the chain-level failure at
each degree, and the total cancellation $\sum_k \delta^{(k)} = 0$
is enforced by the TCFT structure. This per-$k$ vs.\ total
distinction is the source of cached confusion \#14: whoever
asserts ``$[m_k, B^{(2)}] = 0$ for each $k$'' has substituted
\emph{strict} Stasheff identities for the genuinely
\emph{homotopy-coherent} TCFT identity, and this substitution is
silently wrong outside the SC-formal locus.
\end{remark}

\begin{remark}[Analogy with spectral $z$ vs worldsheet $z$ (AP-CY31)]
\label{rem:m3b2-category-error}
The chain-level $[m_3, B^{(2)}] \neq 0$ is a \emph{category
error} of the same type as AP-CY31 (spectral $z$ vs
worldsheet~$z$).  In both cases:
\begin{itemize}
 \item A quantity is computed in one framework (strict chain
  level, resp.\ worldsheet coordinates).
 \item The nonvanishing is real \emph{in that framework}.
 \item The correct mathematical question lives in a different
  framework (homotopy-coherent/$\infty$-categorical, resp.\
  spectral/Yangian).
 \item Transfer to the correct framework resolves the apparent
  obstruction.
\end{itemize}
The general lesson: when a computation produces a nonzero
obstruction, verify that the computation lives in the same
category as the question before concluding that the programme
fails.
\end{remark}


% ============================================================
\section{Summary of computational engines}
\label{sec:m3b2-engines}
% ============================================================

The following engines, totalling $500+$ tests, were involved in
resolving the $[m_3, B^{(2)}]$ question.  Each engine is an
independent module in \texttt{compute/lib/} with a corresponding
test suite in \texttt{compute/tests/}.

\begin{center}
\renewcommand{\arraystretch}{1.15}
\small
\begin{tabular}{lrl}
 \toprule
 \textbf{Engine} & \textbf{Tests} & \textbf{Role} \\
 \midrule
 \texttt{obs\_ainf\_local\_p2.py} & 54 &
 Explicit $[m_3, B^{(2)}]$ for local $\bP^2$ \\
 \texttt{obs\_ainf\_random\_search.py} & 67 &
 Random search over cyclic $\Ainf$-algebras \\
 \texttt{obs\_ainf\_counterexample\_search.py} & 49 &
 Adversarial search for vanishing \\
 \texttt{ks\_cyclic\_minimal\_obs\_ainf.py} & -- &
 KS minimal model, $117/1280$ nonzero \\
 \texttt{connes\_b\_obs\_ainf.py} & -- &
 Bidegree decomposition analysis \\
 \texttt{spectral\_seq\_obs\_ainf.py} & -- &
 Spectral sequence, non-adjacent terms \\
 \texttt{tsygan\_formality\_obs\_ainf.py} & -- &
 Tsygan formality scope \\
 \texttt{stasheff\_cancellation\_obs\_ainf.py} & 40 &
 Weight-$s$ identities, cross-hierarchy \\
 \texttt{operadic\_tcft\_mk\_b2\_engine.py} & 43 &
 TCFT operadic proof structure \\
 \texttt{hopf\_fibration\_s3\_framing.py} & 67 &
 $S^1 \hookrightarrow S^3 \twoheadrightarrow S^2$
 fibration decomposition of the $S^3$-framing
 (not the AP-CY72 $\mathbb{S}^4 = \mathbb{S}^2 \times
 \mathbb{S}^2$ shorthand) \\
 \texttt{derived\_framing\_obstruction.py} & 51 &
 $\HH^{-2}_{E_1}$ vanishing, Goodwillie \\
 \texttt{formality\_ainf\_dcoh.py} & -- &
 Formality hierarchy for CY manifolds \\
 \texttt{cya3\_stress\_test.py} & 45 &
 Adversarial stress test of all proofs \\
 \bottomrule
\end{tabular}
\end{center}

\begin{remark}[Computational methodology]
\label{rem:computational-methodology}
The engines use exact rational arithmetic
(\texttt{fractions.Fraction}) throughout, avoiding floating-point
errors that could mask or create spurious nonvanishing.  Each
engine verifies its own input data: Stasheff identities, cyclic
invariance, CY pairing nondegeneracy, and grading consistency
are checked before any commutator computation.  The test suites
are integrated into the Vol~III \texttt{make test} infrastructure
($19{,}838$~total tests).
\end{remark}
